\subsection{单步准确率不足以评价长期预测能力}

单步预测的输入始终来自真实历史，而规划和想象轨迹会递归使用模型自己的预测。M1000 中多个模型的单步 NRMSE 很接近，但 H20 和 H50 差距明显；单步误差最低的模型也不一定具有最低的 H5--H20 综合误差。该结果表明，面向控制的世界模型需要联合评价单步预测和多步递归预测，不能只依据单步损失决定部署模型。

\subsection{预测精度与控制效果并不完全一致}

M1000 控制变量实验中，统一预测评价选择 Transformer-2L；MPPI 在 Transformer-4L 上取得更高的策略比较阶段回报，MB-PPO 则在 LSTM 上表现更好。这一现象提示，通用预测误差与下游控制效果可能并不完全一致，不同策略也可能对不同类型的模型误差具有不同敏感性。该结论仅是当前五种架构、两类策略和一个数据规模上的实验观察，不应上升为普遍规律，也不用于反向修改世界模型选择。

\subsection{世界模型误差可能被策略优化放大}

M100 的结果显示，更复杂或采样更积极的优化器并不一定更可靠。规划器可能找到世界模型预测回报很高、但真实仿真表现较差的动作区域；当训练数据覆盖不足时，这类动作可能位于模型不可信的区域。因此，基于模型的控制不仅要优化预测回报，还需要判断候选动作是否受到历史数据支持。

\subsection{数据增加不一定带来单调的控制收益}

从 M100 到 M1000，最终累计奖励明显提高；从 M1000 到 M10000 则没有继续单调提高，M1000+MPPI 的最终平均回报略高。该结果可能与数据覆盖范围、模型架构、训练预算结构以及规划器对局部误差的敏感性共同有关。现有实验只能说明更多数据没有自动转化为更高控制回报，不能把差异唯一归因于某一个因素。

\subsection{模型架构与数据规模存在交互}

M100 中 GRU 取得更低的综合预测误差，提示递归结构在当前小数据设置中具有一定优势；M1000 和 M10000 中轻量的 Transformer-2L 更优；Transformer-4L 没有稳定获得进一步收益。因此，架构选择需要结合数据规模、递归预测长度和后续策略共同考虑，而不应预设 Transformer 或循环网络必然获胜。
